%! AP Statistics — Unit 4, Lesson 4.11 — Student Packet Cover Sheet
\documentclass[10pt]{article}
\usepackage{apstats-article}
\usepackage{apstats-boxes}
\usepackage{ltablex}
\keepXColumns

\begin{document}

% Full-bleed navy banner
\begin{tikzpicture}[remember picture, overlay]
  \fill[navy] (current page.north west)
    rectangle ([yshift=-0.9in]current page.north east);
\end{tikzpicture}
\vspace{-0.8in}

\begin{center}
  {\color{white}\textbf{\LARGE Statistics}}\\[4pt]
  {\color{skymid}\large\textbf{Unit 4: Probability, Random Variables, and Probability Distributions}}\\[6pt]
  {\color{white}\textbf{\large Lesson 4.11 \quad Parameters for a Binomial Distribution}}
\end{center}

\vspace{0.22in}
\namedateperiod
\vspace{0.18in}

\begin{learningtargetbox}
\small
\begin{enumerate}[leftmargin=1.5em, itemsep=5pt]
  \item \ldots{} calculate the mean $\mu_X = np$ and standard deviation
        $\sigma_X = \sqrt{np(1-p)}$ of a binomial distribution, and interpret
        $\mu_X$ as the long-run average number of successes.
  \item \ldots{} compute cumulative binomial probabilities $P(X \le k)$ and
        $P(X \ge k)$ using \texttt{binomcdf} and the complement rule.
\end{enumerate}
\end{learningtargetbox}

\vspace{0.14in}

\begin{tocbox}
% A CONTENTS LIST ONLY: no Score column and no Total row. Whether a component counts is said in
% its description.
\small
\renewcommand{\arraystretch}{1.5}
\begin{tabularx}{\linewidth}{@{} c l X @{}}
\rowcolor{sky}
\textbf{\#} & \textbf{Component} & \textbf{Description} \\
\midrule
1 & Warm-Up        & Spiral review: binomial formula, mean/SD of discrete RV \\
2 & Guided Notes   & $\mu_X = np$; $\sigma_X$; \texttt{binomcdf}; interpretation \\
3 & Group Activity & Computing and interpreting binomial parameters \\
4 & Exit Ticket    & Defective products: $\mu_X$, $\sigma_X$, cumulative $P$ \\
5 & Homework       & Four scenarios; AP free-response \\
\end{tabularx}
\end{tocbox}

\vspace{0.14in}

\begin{remindbox}
\small
The \textbf{most important ideas in Lesson 4.11}: use the formulas, show the setup,
and always interpret $\mu_X$ in context.

\vspace{6pt}
\begin{tabularx}{\linewidth}{@{} >{\centering\arraybackslash\columncolor{goldbg}}p{0.06\linewidth}
                                    X @{}}
\textbf{\large$\mu$} &
\textbf{Mean:} $\mu_X = np$.\enspace
Interpretation: ``If this binomial experiment were repeated many times,
the average number of successes would be $np$.'' \\[6pt]
\textbf{\large$\sigma$} &
\textbf{Standard deviation:} $\sigma_X = \sqrt{np(1-p)}$.\enspace
Don't forget the square root! \\[6pt]
\textbf{\large cdf} &
\textbf{Cumulative:} \texttt{binomcdf(n, p, k)} $= P(X \le k)$.\enspace
For $P(X \ge k)$, use $1 - \texttt{binomcdf(n, p, k-1)}$ --- note $k-1$,
not $k$. \\
\end{tabularx}
\end{remindbox}

\end{document}
